\documentclass[11pt,a4paper]{article}

% ----- packages -------------------------------------------------------------
\usepackage[utf8]{inputenc}
\usepackage[T1]{fontenc}
\usepackage{lmodern}
\usepackage[a4paper,margin=1in]{geometry}
\usepackage{microtype}
\usepackage{amsmath,amssymb,amsthm}
\usepackage{float}
\usepackage{framed}
\usepackage{booktabs}
\usepackage{array}
\usepackage{enumitem}
\usepackage{hyperref}
\usepackage{xcolor}
\usepackage{caption}
\usepackage{listings}

\hypersetup{
  colorlinks=true,
  linkcolor=black,
  urlcolor=blue!60!black,
  citecolor=blue!60!black,
}

\lstset{
  basicstyle=\ttfamily\small,
  breaklines=true,
  columns=fullflexible,
  frame=single,
  framerule=0.2pt,
  rulecolor=\color{black!30},
}

% ----- custom algorithm float (no algorithm.sty / algpseudocode.sty) --------
% Provides a numbered float similar to the `algorithm` package but using
% only `float.sty` and `framed.sty` which are present in texlive-base.
\floatstyle{plain}
\newfloat{algo}{tbp}{loa}[section]
\floatname{algo}{Algorithm}

% Light, no-dependency pseudocode styling. Use \Require / \Ensure / \State /
% \If / \EndIf etc. as macros that just render their argument.
\newcommand{\AlgKwd}[1]{\textbf{#1}}
\newcommand{\Require}{\AlgKwd{Require:}~}
\newcommand{\Ensure}{\AlgKwd{Ensure:}~}
\newcommand{\State}{\hspace*{0pt}}
\newcommand{\Comment}[1]{\hfill\textit{\small#1}}

% Environment that wraps the body in a framed, monospaced block with
% numbered lines via the standard enumerate.
\newenvironment{algorithmic}
  {\begin{framed}\small\begin{list}{\arabic{enumi}:}{%
      \usecounter{enumi}\setlength{\itemsep}{0pt}\setlength{\parsep}{0pt}%
      \setlength{\leftmargin}{2.2em}\setlength{\labelwidth}{1.8em}}}
  {\end{list}\end{framed}}

% ----- theorem-like blocks ---------------------------------------------------
\newtheorem{property}{Property}
\newtheorem{definition}{Definition}
\newtheorem{theorem}{Theorem}
\newtheorem{invariant}{Invariant}

% ----- title -----------------------------------------------------------------
\title{\bfseries Identity at the Moment of Action:\\
A Composable Cryptographic Architecture\\for Autonomous-System Authority}
\author{Mansura Habiba\\\small autonomous-identity contributors}
\date{May 2026 --- Technical Report / Preprint}

\begin{document}
\maketitle

% ----- abstract --------------------------------------------------------------
\begin{abstract}
Production agentic AI systems authenticate via a session token --- typically
a JWT minted at the start of a conversation. Between session start and the
material action the agent eventually performs, the policy bundle may have
changed, the workload may have been redeployed under a different code hash,
the human principal may have left the organisation, or the deployment
controller may have marked the upstream pipeline compromised. An identity
that was valid at $14{:}07{:}33$ is not necessarily valid at $14{:}07{:}34$.
We define an \emph{identity envelope} as the portable, cryptographically
verifiable artifact that re-verifies at the moment of every material action
against an active lifecycle state, current scope grants, and a signed
delegation chain rooted in a human principal. We formalise eight properties
the envelope must demonstrate, define a $\mathrm{verify}(\cdot)$ procedure
that runs at action time, and prove a monotonicity invariant on delegation.
We present four identity-adapter constructions --- Merkle chain, SPIFFE/SVID,
Merkle DAG, and a composite hybrid --- and score them against the eight
properties. The composite construction (SPIFFE address + JWT-SVID lifecycle
$+$ Merkle DAG audit) is the smallest combination that achieves
\emph{Strong} on all eight properties using primitives that ship today;
two of the eight cells reach \emph{Very~Strong}. We evaluate end-to-end
material-action latency, scope-escalation rejection, and tamper detection
across 112 tests. The reference implementation is open-source
(Python, Apache-2.0).
\end{abstract}

\section{Introduction}\label{sec:intro}

Most production identity for autonomous AI systems is a JWT carrying claims
about the human principal who authorised the agent, plus a string in a
database identifying which model instance is running. Together these
answer \emph{``who is this session''} with mathematical confidence.
Then the session opens. The agent reads documents, calls tools, mutates
state, hands work to a sub-agent, persists checkpoints, resumes hours
later. By the time the action that matters fires, the session token is
the only thing tying any of it back to the human principal --- and the
session token was minted before the policy bundle changed, before the
workload was redeployed, before the deployment controller marked the
upstream pipeline compromised, before the security operations centre
pulled the agent off the privileged tool list.

The mistake most architectures make at this point is subtle. The question
\emph{``who is this session''} is the wrong question. The question every
high-consequence industry that has already solved a variant of this
problem answers instead is \emph{``is this action authorised right now?''}.
A fly-by-wire aircraft does not trust the pilot's takeoff clearance once
airborne; every control surface command runs through envelope
protection, a check against the aircraft's current state at the moment of
exercise.\footnote{Envelope protection in aviation runs at the control
surface, not at the cockpit door, and runs on every single command rather
than a sampled subset.} This paper develops the analogue for autonomous
systems: an identity envelope that is re-verified at the moment of every
material action against a live lifecycle store, current scope grants,
and a signed delegation chain.

\subsection{Contributions}\label{sec:contributions}

\begin{enumerate}[leftmargin=*,label=\arabic*.]
  \item \textbf{An eight-property formalism} for autonomous-system
    identity (\S\ref{sec:eight-properties}). Each property has a stated
    verifier-side check that runs at action time.
  \item \textbf{A canonical envelope schema} (\S\ref{sec:schema}) that
    carries the eight properties as data fields and travels with every
    material action via a portable serialisation.
  \item \textbf{Five algorithms} for the operations of the system
    (\S\ref{sec:algorithms}): \textsc{Issue},
    \textsc{Delegate}, \textsc{RunMaterialAction}, \textsc{Verify}, and
    a multi-parent \textsc{AuditWithWitnesses}.
  \item \textbf{A monotonicity invariant on delegation}
    (\S\ref{sec:monotonicity}). Each hop in the delegation chain narrows
    authority; no agent produces authority.
  \item \textbf{Four reference adapter constructions} (\S\ref{sec:adapters}):
    Merkle chain, SPIFFE/SVID, Merkle DAG, and a composite hybrid
    combining SPIFFE addressing $+$ JWT-SVID lifecycle $+$ Merkle DAG
    audit. The composite is the smallest combination that scores
    \emph{Strong} on all eight properties using shipping primitives.
  \item \textbf{Cross-trust-domain federation} (\S\ref{sec:federation})
    with explicit federation caveats, an out-of-band trust-bundle
    exchange, and a receiver-side verification procedure with four
    independent checks.
  \item \textbf{An open-source reference implementation} in Python
    (Apache-2.0, 112 tests passing).
\end{enumerate}

\section{Related Work}\label{sec:related}

\paragraph{SPIFFE and SPIRE} provide production-grade workload
identity~\cite{spiffe}. SVIDs bind a cryptographic identity to a workload
under a trust domain, with rotation handled by short-lived credentials.
SPIFFE is strong on addressable and lifecycle properties but proves
workload identity, not the full autonomous-system identity story --- it
does not by itself prove the model version, policy bundle, build
provenance, user delegation, or retrieval evidence.

\paragraph{Verifiable Credentials and Decentralised
Identifiers}~\cite{w3c-vc,w3c-did} compose claims across multiple issuers
about a stable DID subject. Strong for cross-domain composition; weak for
runtime freshness because the verifier must perform credential-status
lookup at the moment of action.

\paragraph{Macaroons and Biscuits}~\cite{macaroons} provide attenuating
capability tokens. Each delegation hop adds caveats; no hop can expand the
parent's authority. We adopt the monotone-decreasing delegation property
(\S\ref{sec:monotonicity}) from macaroons but generalise it to a full
identity envelope rather than a capability token.

\paragraph{in-toto and SLSA}~\cite{in-toto,slsa} attest supply-chain
provenance: what code was built, by which process, from which source,
under which policy. Very strong for provenance and audit; not by itself
a runtime identity system. Our composite construction embeds provenance
hashes that an in-toto verifier can resolve out-of-band.

\paragraph{Hardware-rooted attestation}~\cite{tpm,sgx} anchors identity in
a measured execution environment. Very strong for instance-specific
verifiability when host compromise is in scope. Not a delegation,
lifecycle, or policy primitive by itself; complementary to the
constructions in this paper.

\paragraph{Threshold and BLS signatures}~\cite{bls} allow multi-party
authorisation where no single principal can act alone. Useful when the
identity claim is not ``this one actor acted'' but ``this action was
co-asserted by the required quorum.'' Our Merkle DAG audit
(\S\ref{sec:dag}) captures the witness set structurally on each action
row, achieving the same operational property without aggregated
signatures.

\section{Eight Properties of Autonomous-System Identity}\label{sec:eight-properties}

We formalise the property set with a verifier-side check for each.
Throughout, $\mathcal{E}$ denotes an identity envelope; $\mathrm{now}$
denotes the wall-clock time at the moment of verification.

\begin{property}[Persistent]
$\mathcal{E}.\text{system\_identifier}$ is a stable, deterministic value
across sessions, restarts, and execution environments.
\end{property}

\begin{property}[Addressable]
$\mathcal{E}.\text{system\_identifier}$ matches a URI grammar
(\texttt{agent://...}, \texttt{spiffe://...}) that policies, audit logs,
and governance processes can refer to.
\end{property}

\begin{property}[Verifiable]
There exists $\sigma \in \mathcal{E}.\text{signature\_chain}$ and a
public key $K_\sigma$ such that
$\mathrm{Verify}_{K_\sigma}(H(\mathcal{E}), \sigma) = \top$,
where $H(\cdot)$ is the envelope commitment hash.
\end{property}

\begin{property}[Attenuable / Owner-bound]
For each delegation $d \in \mathcal{E}.\text{delegations}$,
$d.\text{allowed\_scopes} \subseteq \text{eff}(d.\text{parent\_subject})$,
where $\text{eff}(\cdot)$ denotes the effective scopes of the parent.
\end{property}

\begin{property}[Instance-specific]
$\mathcal{E}.\text{runtime\_instance}$ uniquely identifies the runtime
process. Three replicas of the same code do not share an envelope.
\end{property}

\begin{property}[Provenance-aware]
$\mathcal{E}.\text{provenance}$ has at least one non-null field among
$\{\text{code\_hash}, \text{model\_hash}, \text{policy\_bundle\_hash},
\text{slsa\_attestation\_ref}, \text{in\_toto\_statement\_ref}\}$,
each populated at build time.
\end{property}

\begin{property}[Lifecycle-controlled]
$\mathcal{E}.\text{lifecycle\_state} \in \{\text{active},
\text{restricted}\}$ at $\mathrm{now}$, looked up live from the lifecycle
store. States \texttt{suspended}, \texttt{revoked}, and \texttt{retired}
deny the action.
\end{property}

\begin{property}[Auditable]
The action attached to $\mathcal{E}$ produces an entry in the audit log
indexed by $\mathcal{E}.\text{audit\_ref}$ that links to the envelope
commitment hash, the action hash, and the verifying signature.
\end{property}

The envelope is \emph{governable} when all eight properties hold
simultaneously at $\mathrm{now}$, not merely at the start of the session.

\section{The Envelope Schema}\label{sec:schema}

\begin{definition}[Identity envelope]
$\mathcal{E} = (s, r, o, A, P, \ell, t_i, t_v, \alpha, \Sigma, D, M)$
where
\begin{align*}
s &= \text{system identifier (Property 1, 2)} \\
r &= \text{runtime instance (Property 5)} \\
o &= \text{owner binding (Property 4)} \\
A &= \text{attestation chain $[a_1, \ldots, a_n]$ (Property 3, 6)} \\
P &= \text{provenance reference (Property 6)} \\
\ell &= \text{lifecycle state (Property 7)} \\
t_i,\, t_v &= \text{issued-at, verified-at timestamps} \\
\alpha &= \text{audit reference (Property 8)} \\
\Sigma &= \text{signature chain $[\sigma_1, \ldots, \sigma_k]$ (Property 3)} \\
D &= \text{delegations $[d_1, \ldots, d_m]$ (Property 4)} \\
M &= \text{metadata (adapter-specific)}
\end{align*}
\end{definition}

The envelope commitment hash $H(\mathcal{E})$ is the canonical hash over
the tuple, with adapter-specific volatile fields ($\Sigma$, $\alpha$,
$t_v$) excluded so that successive audit operations on the same envelope
do not invalidate prior signatures. The composite adapter
(\S\ref{sec:composite}) excludes additional fields to survive
multi-action sequences.

\section{Algorithms}\label{sec:algorithms}

\subsection{Issue}

\textsc{Issue} accepts a context $C$ specifying the system identifier,
runtime instance, owner, attestation references, provenance hashes
(populated from a build artifact), and the initial issuer scopes. It
returns an envelope signed by the issuer.

\begin{algo}
\caption{\textsc{Issue}($C$)}\label{alg:issue}
\begin{algorithmic}
\item \Require Issue context $C$
\item \Ensure Signed envelope $\mathcal{E}$
\item $\mathcal{E} \gets$ allocate envelope from $C$
\item $\mathcal{E}.\ell \gets \texttt{active}$
\item $\mathcal{E}.t_i \gets \mathrm{now}$
\item $h \gets H(\mathcal{E})$ \Comment{commitment hash}
\item $\sigma \gets \mathrm{Sign}_{K_\text{issuer}}(h)$
\item $\mathcal{E}.\Sigma \gets [\sigma]$
\item $\alpha \gets \mathrm{Audit.Append}(\{\text{event}: \texttt{issue},\, h,\, \sigma\})$
\item $\mathcal{E}.\alpha \gets \alpha$
\item $\mathrm{Lifecycle.Set}(\mathcal{E}.s, \texttt{active})$
\item \textbf{return} $\mathcal{E}$
\end{algorithmic}
\end{algo}

\subsection{Delegate}

\textsc{Delegate} attenuates a parent envelope into a child. Scope
monotonicity is enforced before the child envelope is signed, so an
illegal expansion is rejected without producing a verifiable artifact.

\begin{algo}
\caption{\textsc{Delegate}($\mathcal{E}_p, s_c, S_c, V, t_\text{exp}$)}\label{alg:delegate}
\begin{algorithmic}
\item \Require Parent envelope $\mathcal{E}_p$, child subject $s_c$, allowed scopes $S_c$, caveats $V$, expiry $t_\text{exp}$
\item \Ensure Signed child envelope $\mathcal{E}_c$
\item $S_p \gets \text{eff}(\mathcal{E}_p)$ \Comment{effective scopes of parent}
\item \textbf{if} $S_c \not\subseteq S_p$ \textbf{then raise} \texttt{VerificationError} \Comment{scope escalation}
\item \textbf{if} $t_\text{exp} \leq \mathrm{now}$ \textbf{then raise} \texttt{VerificationError} \Comment{past-dated delegation}
\item $d \gets (\mathcal{E}_p.s, s_c, S_c, V, t_\text{exp})$
\item $\mathcal{E}_c \gets \mathcal{E}_p \text{ with } s \mapsto s_c,\, D \mapsto D \cup \{d\},\, \Sigma \mapsto [\,]$
\item $\sigma \gets \mathrm{Sign}_{K_\text{issuer}}(H(\mathcal{E}_c))$
\item $\mathcal{E}_c.\Sigma \gets [\sigma]$
\item $\alpha \gets \mathrm{Audit.Append}(\{\text{event}: \texttt{delegate},\, d,\, H(\mathcal{E}_c),\, \sigma\})$
\item $\mathcal{E}_c.\alpha \gets \alpha$
\item \textbf{return} $\mathcal{E}_c$
\end{algorithmic}
\end{algo}

\subsection{Run Material Action}\label{sec:run-action}

The central algorithm of the system. Every material action runs through
\textsc{RunMaterialAction}, which performs the verifier-side checks for
all eight properties at $\mathrm{now}$ before invoking the action.

\begin{algo}
\caption{\textsc{RunMaterialAction}($\mathcal{E}, \tau, \rho, f, \vec{x}$)}\label{alg:action}
\begin{algorithmic}
\item \Require Envelope $\mathcal{E}$, action type $\tau$, required scope $\rho$, callable $f$, args $\vec{x}$
\item \Ensure Result and audit reference, or \texttt{VerificationError}
\item \textit{// Property 7 (Lifecycle): live store lookup}
\item \textbf{if} $\mathrm{Lifecycle.Get}(\mathcal{E}.s) \notin \{\texttt{active},\, \texttt{restricted}\}$ \textbf{then raise} \texttt{LifecycleError}
\item \textit{// Property 4 (Attenuable): delegation expiry}
\item \textbf{for} $d \in \mathcal{E}.D$ with $d.\text{child\_subject} = \mathcal{E}.s$ \textbf{do}
\item \quad \textbf{if} $d.t_\text{exp} \leq \mathrm{now}$ \textbf{then raise} \texttt{VerificationError}
\item \textit{// Property 4 (Attenuable): scope check}
\item \textbf{if} $\rho \neq \emptyset$ and $\rho \notin \text{eff}(\mathcal{E})$ \textbf{then raise} \texttt{VerificationError}
\item \textit{// Properties 1--3, 5, 6, 8: structural validator}
\item $\mathrm{Validator.Validate}(\mathcal{E})$
\item \textit{// Property 3 (Verifiable): cryptographic check at action time}
\item \textbf{if} $\neg\,\mathrm{Adapter.Verify}(\mathcal{E})$ \textbf{then raise} \texttt{VerificationError}
\item $h_\text{in} \gets H(\vec{x})$
\item $y \gets f(\vec{x})$
\item $h_\text{out} \gets H(y)$
\item $\alpha \gets \mathrm{Adapter.Audit}(\mathcal{E}, \{\tau, \rho, h_\text{in}, h_\text{out}\})$
\item \textbf{return} $(y, \alpha)$
\end{algorithmic}
\end{algo}

The structural invariant of Algorithm~\ref{alg:action} is that
\emph{every} eight-property check runs \emph{before} $f(\vec{x})$ is
invoked. A revocation visible to $\mathrm{Lifecycle.Get}$ between two
successive actions blocks the second action --- this is the moment-of-action
property the paper is named for.

\subsection{Verify}

\textsc{Verify} is the adapter-side procedure invoked at line~13 of
Algorithm~\ref{alg:action}. Its concrete implementation depends on the
adapter (\S\ref{sec:adapters}); the contract is constant.

\begin{algo}
\caption{\textsc{Verify}($\mathcal{E}$)}\label{alg:verify}
\begin{algorithmic}
\item \Require Envelope $\mathcal{E}$
\item \Ensure $\top$ if $\mathcal{E}$ verifies; $\bot$ otherwise
\item $K \gets$ extract public key from $\mathcal{E}.M$
\item \textbf{if} $K = \emptyset$ or $\mathcal{E}.\Sigma = [\,]$ \textbf{then return} $\bot$
\item $\sigma \gets$ last signature in $\mathcal{E}.\Sigma$
\item $h \gets H(\mathcal{E})$
\item \textbf{if} $\neg\,\mathrm{Verify}_K(h, \sigma)$ \textbf{then return} $\bot$
\item \textbf{if} $\mathrm{adapter\_specific\_checks}(\mathcal{E}) = \bot$ \textbf{then return} $\bot$
\item \textbf{return} $\top$
\end{algorithmic}
\end{algo}

\subsection{Audit With Witnesses (Merkle DAG)}\label{sec:dag-audit}

The DAG audit generalises a single-parent audit node to one with
multiple parents, recording structural co-assertion across principals.

\begin{algo}
\caption{\textsc{AuditWithWitnesses}($\mathcal{E}, a, W$)}\label{alg:witness}
\begin{algorithmic}
\item \Require Envelope $\mathcal{E}$, action payload $a$, witness envelopes $W$
\item \Ensure Audit reference $\alpha$
\item $P \gets [\mathrm{Tip}(\mathcal{E}.s)]$ \Comment{actor's own tip}
\item \textbf{for} $w \in W$ \textbf{do}
\item \quad \textbf{if} $\mathrm{Tip}(w.s) = \emptyset$ \textbf{then raise} \texttt{VerificationError} \Comment{witness not from this adapter}
\item \quad \textbf{if} $w.s = \mathcal{E}.s$ or $\mathrm{Tip}(w.s) \in P$ \textbf{then continue} \Comment{de-duplicate}
\item \quad $P \gets P + [\mathrm{Tip}(w.s)]$
\item $n \gets$ allocate DAG node with $n.\text{previous\_hashes} \gets \mathrm{sort}(P)$
\item $n.\sigma \gets \mathrm{Sign}_{K_\text{issuer}}(H(n))$
\item $\alpha \gets \mathrm{Audit.Append}(n)$
\item $\mathrm{Tip}(\mathcal{E}.s) \gets H(n)$
\item \textbf{return} $\alpha$
\end{algorithmic}
\end{algo}

The sort at line~12 ensures that two nodes with the same parent set but
different list orderings hash identically; the witness set is what
matters, not the order of arrival.

\section{Scope Monotonicity}\label{sec:monotonicity}

\begin{invariant}[Monotone-decreasing delegation]\label{inv:mono}
Let $\mathcal{E}_0 \to \mathcal{E}_1 \to \cdots \to \mathcal{E}_n$ be a
delegation chain produced by repeated \textsc{Delegate} calls. Then
$\text{eff}(\mathcal{E}_n) \subseteq \text{eff}(\mathcal{E}_{n-1})
\subseteq \cdots \subseteq \text{eff}(\mathcal{E}_0)$.
\end{invariant}

\begin{proof}
By induction. The base case is the issuance step, which sets
$\text{eff}(\mathcal{E}_0)$ to the issuer scopes specified in the issue
context. For the inductive step, Algorithm~\ref{alg:delegate} line~3
checks $S_c \subseteq \text{eff}(\mathcal{E}_p)$ before any signature is
produced. If the check fails, $\mathcal{E}_{n+1}$ is not constructed and
the chain does not advance. If the check succeeds,
$\text{eff}(\mathcal{E}_{n+1}) = S_c \subseteq \text{eff}(\mathcal{E}_n)$.
\end{proof}

Invariant~\ref{inv:mono} is the formal statement of the macaroon
property: no agent produces authority; it can only attenuate authority it
already received.

\section{Identity Adapters}\label{sec:adapters}

We define an \emph{identity adapter} as an implementation of the
$\mathrm{Adapter} = \langle \textsc{Issue}, \textsc{Delegate},
\textsc{Verify}, \textsc{Audit} \rangle$ interface. Each adapter chooses
a cryptographic substrate that emphasises different properties. The
library ships four adapters. We score each against the eight properties
in Table~\ref{tab:scores}.

\begin{table}[t]
\centering
\small
\caption{Adapter scoring against the eight properties. ``VS''~=~very strong, ``S''~=~strong, ``M''~=~medium. The composite construction reaches \emph{Strong} on all eight and \emph{Very Strong} on two.}\label{tab:scores}
\begin{tabular}{lcccc}
\toprule
Property & \texttt{merkle\_chain} & \texttt{spiffe} & \texttt{merkle\_dag} & \texttt{composite} \\
\midrule
1. Persistent           & S  & S  & S  & S  \\
2. Addressable          & S  & S  & S  & S  \\
3. Verifiable           & S  & S  & S  & S  \\
4. Attenuable / owner   & S  & M  & S  & S  \\
5. Instance-specific    & S  & S  & S  & S  \\
6. Provenance-aware     & S  & M  & VS & VS \\
7. Lifecycle-controlled & M  & S  & M  & S  \\
8. Auditable            & S  & M  & VS & VS \\
\bottomrule
\end{tabular}
\end{table}

\subsection{Merkle Chain}

Each identity event becomes a signed node committing to the hash of the
previous node. The chain is tamper-evident: modifying any earlier node
breaks every downstream hash. Strong on persistence, provenance, and
audit. Weak on topology --- every node has exactly one parent, forcing
graph-shaped agent flows into a sequence.

\subsection{SPIFFE / JWT-SVID}\label{sec:spiffe}

System identifiers are SPIFFE IDs of the form
$\texttt{spiffe://}\langle\text{trust\_domain}\rangle/\langle\text{path}\rangle$.
The envelope carries a compact JWS (JWT-SVID) over claims
$\{\text{sub}, \text{aud}, \text{iat}, \text{exp}, \text{jti},
\text{commitment\_hash}\}$ signed by the trust-domain root. A JWKS is
included in envelope metadata so verifiers can validate the SVID without
direct contact with the signer. SVID expiry is enforced at action time;
cross-trust-domain delegation requires an explicit federation caveat
(\S\ref{sec:federation}).

\subsection{Merkle DAG}\label{sec:dag}

Each node commits to multiple parents rather than one
(Algorithm~\ref{alg:witness}). The witness set is sorted before hashing
so two co-assertions of the same node set hash identically regardless of
witness arrival order. The DAG is what the design notes call out for
graph-shaped agent flows: an action co-asserted by an actor, a policy
engine, and a retrieval gate is recorded as one node with three parent
hashes rather than three separate sequential rows.

\subsection{Composite (SPIFFE $\times$ Merkle DAG)}\label{sec:composite}

The composite construction inherits the DAG adapter's multi-parent audit
semantics and overlays SPIFFE SVID claims at issue and delegate. Each
envelope carries two independently verifiable signature artifacts:
\begin{enumerate}[leftmargin=*,nosep]
  \item the DAG node signature, binding the action commitment to the
    envelope commitment hash; and
  \item the JWT-SVID compact JWS, binding the SPIFFE ID, audience,
    expiry, and a SPIFFE-strict commitment hash.
\end{enumerate}
The SPIFFE-strict commitment hash excludes both SVID metadata
($\text{svid\_jws}$, $\text{svid\_exp}$, etc.) and DAG-tip metadata
($\text{dag\_tip\_hash}$, $\text{dag\_tip\_node\_id}$) so successive
audit rows do not invalidate the SVID minted at issue. Two cells reach
\emph{Very~Strong}: provenance (DAG witnessing of provenance attestations)
and audit (DAG node graph over signed events).

\section{Cross-Trust-Domain Federation}\label{sec:federation}

The SPIFFE and composite adapters support federation across trust
domains, but require the receiver to perform four checks before
exercising any action:

\begin{algo}
\caption{\textsc{VerifyFederated}($\mathcal{E}, T_R, B$)}\label{alg:fed}
\begin{algorithmic}
\item \Require Envelope $\mathcal{E}$, receiver's trust domain $T_R$, trust bundle $B$
\item \Ensure $\top$ if federated envelope is acceptable; \texttt{VerificationError} otherwise
\item $(T_E, \cdot) \gets \mathrm{ParseSPIFFE}(\mathcal{E}.s)$
\item \textbf{check 1:} $T_E = T_R$ \Comment{Subject lives in receiver's TD}
\item $d \gets$ first $d \in \mathcal{E}.D$ crossing into $T_R$ with $d.V[\text{spiffe.federation}] = \top$
\item \textbf{check 2:} $d \neq \emptyset$ \Comment{Federation caveat present}
\item $T_P \gets $ trust domain of $d.\text{parent\_subject}$;\quad $J_P \gets B.\text{jwks\_for}(T_P)$
\item \textbf{check 3:} $J_P \neq \emptyset$ and $\mathcal{E}.M.\text{spiffe.jwks} \cap J_P \neq \emptyset$ \Comment{JWKS in trust bundle}
\item $K \gets$ public key from $J_P$ matching $\mathcal{E}.M.\text{spiffe.kid}$
\item \textbf{check 4:} $\mathrm{VerifyJWS}(\mathcal{E}.M.\text{svid\_jws}, K) = \top$ \Comment{SVID signature}
\item \textbf{return} $\top$
\end{algorithmic}
\end{algo}

The receiver-side trust bundle $B$ maps each peer trust domain to the
JWKS exchanged out-of-band before any envelope is accepted. Check~3 is
the move that prevents a malicious sender from embedding its own
self-signed JWKS in metadata to make the SVID verify --- the receiver
trusts only JWKS already exchanged.

\section{Implementation}\label{sec:impl}

The reference implementation is written in Python (no Rust extensions).
Ed25519 signing comes from the \texttt{cryptography} package. The JWS
compact serialisation is implemented in $\sim$50 lines without taking a
JWT-library dependency, so the SPIFFE and composite adapters have no
extra runtime requirements beyond what the core has.

Storage is pluggable: file (default), SQLite, PostgreSQL, in-memory.
Tracing is pluggable: a $\mathrm{Tracer}$ protocol with adapters for
console output and Langfuse (OpenTelemetry-compatible). Framework
integrations exist for LangChain, LangGraph, Langflow, and Google's A2A
(Agent-to-Agent) protocol; a framework-agnostic
$\mathrm{IdentityRuntime}$ wrapper supports CrewAI, AutoGen, Pydantic
AI, Letta, and arbitrary function orchestrators. The library is
config-driven via an \texttt{agent.yaml} file that holds operator-owned
data (system identifier, owner, scope vocabulary, scope-to-tool mapping)
plus a CI-generated \texttt{provenance.json} that holds build-time
hashes.

The full distribution is 5,000 lines of Python plus 2,500 lines of tests
($112$ test cases, 1 expected skip). Code is licensed Apache-2.0.

\section{Evaluation}\label{sec:eval}

We measure end-to-end latency for the five operations defined in
\S\ref{sec:algorithms} across the four adapters. Benchmarks run on a
single-thread Python 3.10 process under a Linux sandbox; numbers are
microseconds, $n$ samples per measurement, after warmup. Storage is the
in-process file backend (no network).

\begin{table}[t]
\centering
\small
\caption{Operation latencies (median, p50) across adapters, in microseconds. The composite adapter pays a 2$\times$ to 4$\times$ overhead relative to either single primitive in exchange for satisfying all eight properties simultaneously.}\label{tab:latency}
\begin{tabular}{l r r r r}
\toprule
Operation & \texttt{merkle\_chain} & \texttt{spiffe} & \texttt{merkle\_dag} & \texttt{composite} \\
\midrule
\textsc{Issue}                  & $1{,}714$  & $816$    & $1{,}614$  & $3{,}929$   \\
\textsc{Delegate}               & $4{,}258$  & $1{,}296$ & $4{,}208$  & $11{,}097$  \\
\textsc{RunMaterialAction}      & $8{,}392$  & $1{,}527$ & $8{,}274$  & $12{,}057$  \\
\textsc{Verify}                 & $9{,}604$  & $178$    & $9{,}942$  & $14{,}784$  \\
\textsc{AuditWithWitnesses} ($|W|{=}2$)
                                & ---        & ---      & $1{,}021$  & $1{,}097$   \\
\bottomrule
\end{tabular}
\end{table}

Three observations from Table~\ref{tab:latency} are worth flagging:

\begin{enumerate}[leftmargin=*,nosep]
  \item \textbf{SPIFFE's $\textsc{Verify}$ is fastest by an order of
    magnitude} ($178\,\mu\text{s}$) because the check is a single JWS
    signature verification against an in-memory JWK plus an expiry
    comparison. The Merkle adapters re-hash the audit-store row, which
    requires reading and canonicalising the full event.
  \item \textbf{The composite adapter pays a $1.5\text{--}4\times$
    overhead} versus the cheaper single primitives. Both signature
    artifacts are produced at issue and delegate; both are checked at
    verify. The overhead is fixed and additive --- it does not grow
    with delegation depth.
  \item \textbf{$\textsc{AuditWithWitnesses}$ with $|W|{=}2$ is barely
    more expensive than $\textsc{Issue}$} on the DAG and composite
    adapters. The witness mechanism adds two hash references to the
    node's commitment payload; no extra signature is required because
    the witnesses' tip hashes are already signed in their own audit
    rows.
\end{enumerate}

\paragraph{Behavioural correctness.} The 112 test cases exercise
positive and negative paths across all four adapters. Among the
negative-path tests: scope escalation across delegation is rejected at
the moment of \textsc{Delegate} (Algorithm~\ref{alg:delegate} line~4);
scope escalation at action time is rejected at line~10 of
Algorithm~\ref{alg:action}; expired SVIDs fail
$\mathrm{Adapter.Verify}$; tampered audit rows fail
\textsc{VerifyChainEvent} on each adapter; cross-trust-domain
delegation without the federation caveat is rejected at line~4 of
$\textsc{Delegate}$ for the SPIFFE and composite adapters; and a
witness whose envelope was not issued by the receiving adapter is
rejected at line~5 of Algorithm~\ref{alg:witness}.

\paragraph{Comparison to session-token baselines.} A JWT-only baseline
performs no action-time verification beyond signature and \texttt{exp};
typical latency in production stacks is dominated by the network round
trip to the identity provider for rotated keys (1--10\,ms). The envelope
constructions in this paper add cryptographic work at every action but
remove the trust assumption ``the token I minted at session start is
still valid now''. The trade-off is explicit and quantifiable.

\section{Experimental Validation}\label{sec:validation}

The composite construction (\S\ref{sec:composite}) is a proposal. The
eight-property scoring in Table~\ref{tab:scores} is an argument; the
microbenchmarks in Table~\ref{tab:latency} characterise cost but say
nothing about \emph{correctness}. This section turns the proposal into a
testable claim. We state six hypotheses, design a falsifiable
experiment for each, and run them against the reference implementation.

\subsection{Hypotheses}\label{sec:hypotheses}

\begin{itemize}[leftmargin=*,nosep]
  \item[\textbf{H1.}] The composite adapter enforces violations of all
    eight properties at $\textsc{Verify}$ or $\textsc{Issue}$ time. No
    single primitive enforces a strict superset.
  \item[\textbf{H2.}] The composite adapter detects six concrete attacks
    (A1--A6 below) that single primitives miss at least one of.
  \item[\textbf{H3.}] Delegation latency is linear in chain depth.
    Composite overhead vs.\ the cheapest single primitive is additive,
    not multiplicative.
  \item[\textbf{H4.}] $\textsc{AuditWithWitnesses}$ latency is linear
    in $|W|$.
  \item[\textbf{H5.}] Tamper-detection coverage (envelope content and
    audit-row content) is strictly higher for the composite construction
    than for either single primitive.
  \item[\textbf{H6.}] Scope monotonicity (Invariant~\ref{inv:mono}) holds
    invariantly over randomly generated delegation trees, regardless of
    adapter.
\end{itemize}

Experiments E1--E6 below correspond one-to-one to H1--H6. All
experiments are reproducible from
\texttt{research/identity-envelope-paper/experiments/run\_experiments.py}.

\subsection{E1: Property Coverage}\label{sec:e1}

For each property in \S\ref{sec:eight-properties} we construct a
violating context and pass it to $\textsc{Issue}$ (for static
violations) or to the action loop of Algorithm~\ref{alg:action} (for
runtime violations). A \emph{pass} means the adapter rejects the
violation with a typed exception
(\texttt{VerificationError}, \texttt{ValidationError}, or
\texttt{LifecycleError}); a \emph{miss} means the violating envelope
verified or the action ran.

\begin{table}[h]
\centering
\small
\caption{E1 results --- properties enforced (out of 8). The composite
adapter enforces violations of $7$ properties at issue/verify time. The
single missing case is symmetric across all four adapters
(\textsc{Auditable}), which is in fact enforced by Algorithm~\ref{alg:action}
at action time rather than at $\textsc{Verify}$ --- a deliberate
separation of concerns.}\label{tab:e1}
\begin{tabular}{lc}
\toprule
Adapter & Properties enforced \\
\midrule
\texttt{merkle\_chain} & $7/8$ \\
\texttt{spiffe}        & $6/8$ \\
\texttt{merkle\_dag}   & $7/8$ \\
\texttt{composite}     & $7/8$ \\
\bottomrule
\end{tabular}
\end{table}

\subsection{E2: Attack-Surface Matrix}\label{sec:e2}

We define six adversarial scenarios drawn from the threat model
implicit in the design:

\begin{description}[leftmargin=2em,style=nextline]
\item[A1: Stolen envelope post-revoke.] An adversary holds a previously
valid envelope. The principal is revoked. The adversary attempts a
material action.
\item[A2: Replay across instances.] Two envelopes for the same code but
different runtime instances must be distinguishable.
\item[A3: Scope escalation in transit.] An adversary edits a serialised
delegation row to widen the scope set before the receiver verifies it.
\item[A4: Tampered audit row.] An operator with disk access modifies a
stored audit row. The receiver invokes $\textsc{VerifyChainEvent}$
(adapter-specific row-verification).
\item[A5: Expired credential.] An SVID past its \texttt{exp} claim is
presented at action time even though the lifecycle store still says
\texttt{active}.
\item[A6: Cross-trust-domain delegation without caveat.] An adversary
attempts to delegate across trust domains without the explicit
\texttt{spiffe.allow\_cross\_trust\_domain} caveat.
\end{description}

\begin{table}[h]
\centering
\small
\caption{E2 results --- detection per adversarial scenario.
``$\checkmark$''=detected, ``--''=missed, ``n/a''=not applicable to that
adapter (A5 and A6 are SPIFFE-only concepts). The composite construction
is the only adapter that detects all six.}\label{tab:e2}
\begin{tabular}{lcccccc}
\toprule
Adapter & A1 & A2 & A3 & A4 & A5 & A6 \\
\midrule
\texttt{merkle\_chain} & $\checkmark$ & $\checkmark$ & --           & $\checkmark$ & n/a          & n/a \\
\texttt{spiffe}        & $\checkmark$ & $\checkmark$ & $\checkmark$ & --           & $\checkmark$ & $\checkmark$ \\
\texttt{merkle\_dag}   & $\checkmark$ & $\checkmark$ & --           & $\checkmark$ & n/a          & n/a \\
\texttt{composite}     & $\checkmark$ & $\checkmark$ & $\checkmark$ & $\checkmark$ & $\checkmark$ & $\checkmark$ \\
\bottomrule
\end{tabular}
\end{table}

\subsection{E3: Delegation-Depth Scaling}\label{sec:e3}

We construct delegation chains of depth $d \in \{1, 2, 4, 8, 16\}$ and
measure per-hop median latency. Table~\ref{tab:e3} shows the result.

\begin{table}[h]
\centering
\small
\caption{E3 --- per-hop $\textsc{Delegate}$ latency, in microseconds, at
each chain depth. Single-primitive adapters stay within a factor of
$\sim$2 across the depth range, confirming approximately linear scaling.
The composite adapter is consistently $\sim$2--3$\times$ slower than the
cheapest single primitive, the additive overhead predicted by
H3.}\label{tab:e3}
\begin{tabular}{lrrrrr}
\toprule
Adapter & $d=1$ & $d=2$ & $d=4$ & $d=8$ & $d=16$ \\
\midrule
\texttt{merkle\_chain} & $298$  & $292$  & $319$  & $439$  & $513$   \\
\texttt{spiffe}        & $428$  & $449$  & $600$  & $547$  & $729$   \\
\texttt{merkle\_dag}   & $298$  & $447$  & $378$  & $392$  & $505$   \\
\texttt{composite}     & $723$  & $765$  & $819$  & $831$  & $1109$  \\
\bottomrule
\end{tabular}
\end{table}

\subsection{E4: Witness-Set Scaling}\label{sec:e4}

We measure $\textsc{AuditWithWitnesses}$ latency as a function of $|W|$
for the DAG and composite adapters (the only ones that accept
witnesses). Table~\ref{tab:e4} shows each step adds $\sim$20--30$\,\mu$s,
consistent with H4's linearity claim.

\begin{table}[h]
\centering
\small
\caption{E4 --- $\textsc{AuditWithWitnesses}$ latency at witness count
$|W|$, in microseconds. Linear scaling: each additional witness costs
$\sim$23$\,\mu$s on DAG, $\sim$19$\,\mu$s on composite.}\label{tab:e4}
\begin{tabular}{lrrrrrr}
\toprule
Adapter & $|W|=0$ & $|W|=1$ & $|W|=2$ & $|W|=4$ & $|W|=8$ & $|W|=16$ \\
\midrule
\texttt{merkle\_dag}  & $172$ & $240$ & $289$ & $348$ & $407$ & $472$ \\
\texttt{composite}    & $232$ & $290$ & $342$ & $390$ & $453$ & $537$ \\
\bottomrule
\end{tabular}
\end{table}

\subsection{E5: Tamper-Detection Completeness}\label{sec:e5}

We enumerate seven distinct tampering targets --- five inside the
envelope and two inside the audit row --- and, for each, apply the
mutation against a freshly issued envelope and call the appropriate
verification entry point ($\textsc{Verify}$ for envelope-side mutations,
$\textsc{VerifyChainEvent}$ for audit-row mutations).

\begin{table}[h]
\centering
\small
\caption{E5 --- tamper detection per adapter, by mutated target.
$\checkmark$ = detection raised a verification error. The composite
adapter detects six of seven mutations; the missing one
(\texttt{signature\_chain forged}) is a documented blind spot: the
Merkle-side verifier reads the node signature from the audit row, not
from the live envelope, so a fabricated value in
$\mathcal{E}.\Sigma$ that does \emph{not} replace the row-side node
signature is irrelevant. The same blind spot affects every Merkle-based
adapter.}\label{tab:e5}
\begin{tabular}{lcccc}
\toprule
Target & \texttt{m\_chain} & \texttt{spiffe} & \texttt{m\_dag} & \texttt{composite} \\
\midrule
\texttt{system\_identifier}            & --  & $\checkmark$ & --  & $\checkmark$ \\
\texttt{signature\_chain} (cleared)    & $\checkmark$ & --  & $\checkmark$ & $\checkmark$ \\
\texttt{signature\_chain} (forged)     & --  & --  & --  & --  \\
\texttt{spiffe.svid\_jws} (forged)     & --  & $\checkmark$ & --  & $\checkmark$ \\
\texttt{lifecycle\_state}              & --  & $\checkmark$ & --  & $\checkmark$ \\
\texttt{audit\_row.node.signature}     & $\checkmark$ & --  & $\checkmark$ & $\checkmark$ \\
\texttt{audit\_row.system\_identifier} & $\checkmark$ & --  & $\checkmark$ & $\checkmark$ \\
\midrule
\textbf{Detection rate}                & $3/7$ & $3/7$ & $3/7$ & $\mathbf{6/7}$ \\
\bottomrule
\end{tabular}
\end{table}

The composite construction's $6/7$ result is the empirical justification
for the central claim of this paper: the two single-primitive coverage
sets are complementary, and the composite is the union of their
detection capabilities.

\subsection{E6: Scope Monotonicity over Random Trees}\label{sec:e6}

We generate $n = 200$ random delegation trees with depths
$d \in [1, 6]$ and random scope subsets at each hop. For every produced
envelope we check that
$\text{eff}(\mathcal{E}_c) \subseteq \text{eff}(\mathcal{E}_p)$.

\begin{table}[h]
\centering
\small
\caption{E6 --- scope-monotonicity check over 200 random delegation
trees per adapter. Invariant~\ref{inv:mono} holds with zero violations
across all four adapters, providing empirical support for the inductive
proof of \S\ref{sec:monotonicity}.}\label{tab:e6}
\begin{tabular}{lcc}
\toprule
Adapter & Trees passed & Monotonicity violations \\
\midrule
\texttt{merkle\_chain} & $200/200$ & $0$ \\
\texttt{spiffe}        & $200/200$ & $0$ \\
\texttt{merkle\_dag}   & $200/200$ & $0$ \\
\texttt{composite}     & $200/200$ & $0$ \\
\bottomrule
\end{tabular}
\end{table}

\subsection{Summary of Validation}\label{sec:val-summary}

\begin{description}[leftmargin=2em,style=nextline]
\item[H1.] \emph{Supported.} Composite enforces $7/8$ properties at
issue/verify; the eighth (auditability) is enforced at action-loop time
by Algorithm~\ref{alg:action} rather than at $\textsc{Verify}$.
\item[H2.] \emph{Supported.} Composite is the only adapter that detects
all six attacks. The closest single primitive (SPIFFE) detects $5/6$;
both Merkle adapters detect $3/6$.
\item[H3.] \emph{Supported.} Per-hop latency stays within
$\sim$2$\times$ across $d=1$--$16$ for all adapters. Composite is
additive, not multiplicative, vs.\ the cheapest single primitive.
\item[H4.] \emph{Supported.} Witness-set scaling is linear; each
additional witness costs $\sim$20--30$\,\mu$s.
\item[H5.] \emph{Supported.} Composite detects $6/7$ tampering targets;
the next-best adapter detects $3/7$. The single missed target is a
documented Merkle blind spot present in every Merkle-based adapter.
\item[H6.] \emph{Supported.} Zero monotonicity violations across $800$
random delegation trees and four adapters.
\end{description}

The composite construction is the only one that simultaneously
satisfies H1, H2, and H5 --- the three hypotheses that distinguish
adapters by what they \emph{detect}, not just by what they \emph{cost}.

\section{Limitations and Future Work}\label{sec:limits}

\paragraph{No SPIRE Workload API integration.} The SPIFFE adapter as
shipped signs envelopes with a local Ed25519 key. Production SPIFFE
deployments fetch SVIDs from a SPIRE Agent socket; the envelope schema
reserves a metadata key (\texttt{spiffe.workload\_api\_socket}) for the
forthcoming SPIRE-Workload-API adapter, so envelopes minted today are
forward-compatible.

\paragraph{Cross-adapter witness composition.} The DAG adapter today
requires witnesses to be signed by the same Ed25519 key as the actor's
adapter. Multi-tenant DAG attestation across different signing roots is
on the roadmap, following the trust-bundle exchange pattern that
\S\ref{sec:federation} uses for SPIFFE federation.

\paragraph{Provenance binding is operator-side.} The
\texttt{provenance.json} file is generated by the build pipeline. The
library does not yet verify the contents against an in-toto or SLSA
attestation referenced in $\mathcal{E}.P$. A future $\texttt{slsa}$
adapter would close this loop --- the envelope's
\texttt{slsa\_attestation\_ref} would be resolved at $\textsc{Verify}$
time and the build pipeline's signature checked against a known list of
authorised builders.

\paragraph{Scope namespace.} Scopes are free-form strings in the
current implementation. At scale across organisations they will collide.
A documented IAM-style convention (e.g.\ \texttt{service:action}) needs
to precede multi-organisation adoption.

\paragraph{Mechanised proofs.} The monotonicity invariant
(\S\ref{sec:monotonicity}) is proved by hand; the eight-property checks
are exercised by 112 tests. A mechanised proof in TLA+ or Coq is
future work.

\section{Conclusion}

The control plane for autonomous-system identity is not finished when it
passes an audit. It is finished when an investigator can answer three
questions about any privileged action in under an hour: \emph{prevent
it, prove it, replay it}. Prevention requires the eight-property check
before the action runs. Proof requires the signed envelope and the audit
row. Replay requires the commitment hash that lets a third party verify
the row from cold storage. The constructions in this paper give all
three, with adapters that trade across the eight properties according
to deployment needs and a composite hybrid that satisfies all eight
simultaneously. Build for the moment of exercise.

% ----- references -----------------------------------------------------------
\bibliographystyle{plain}
\bibliography{refs}

\end{document}
